\documentclass[10pt,a4paper]{amsart}
\usepackage[margin=19mm]{geometry}
\usepackage{amsmath,amssymb,mathtools}
\usepackage[scr=rsfs,cal=euler]{mathalfa}
\usepackage{xfrac}
\usepackage[unicode,hidelinks,bookmarks=false]{hyperref}
\newcommand{\ie}{i.e.\ }
\newcommand{\cf}{cf.\ }
\newcommand{\resp}{respectively\ }
\newcommand{\Subsection}{\subsection}
\providecommand{\nobreakdash}{\nobreak\hbox{-}}
\setlength{\emergencystretch}{2em}
\multlinegap0pt
\usepackage{bm}
\usepackage{bbm}
\usepackage{dsfont}
\usepackage{xspace}
\usepackage{cancel}
\usepackage{mathtools}
\usepackage{amsthm}
\makeatletter
\@ifundefined{theorem}{\newtheorem{theorem}{Theorem}}{}
\@ifundefined{lemma}{\newtheorem{lemma}[theorem]{Lemma}}{}
\makeatother
\def\Pr{\mathbb{P}}
\def\d{\delta} \def\D{\Delta}
\def\m{\mu} \def\n{\nu} \def\p{\pi}
\def\t{\tau} \def\om{\omega} \def\OM{\Omega} \def\Om{\Omega}\def\up{\upsilon}\def\U{\Upsilon}
\newcommand{\bfrac}[2]{\left(\frac{#1}{#2}\right)}
\newcommand{\brac}[1]{\left(#1\right)}
\newcommand{\wh}[1]{\widehat{#1}}
\title[Cover time of random subgraphs of the hypercube]{Cover time of random subgraphs of the hypercube}
\author{Colin Cooper, Alan Frieze, Wesley Pegden}
\date{Source: arXiv:2506.03375v1 (CC BY 4.0)}
\begin{document}
\maketitle

\noindent\emph{Source and licence.} Cover time of random subgraphs of the hypercube. Version arXiv:2506.03375v1 (CC BY 4.0), distributed under the Creative Commons Attribution 4.0 International licence, as recorded in the arXiv licence metadata for this exact version and shown on the abstract page https://arxiv.org/abs/2506.03375v1. Full rights notice: licenses/arxiv-2506.03375-CC-BY-4.0.txt.\par\smallskip

\noindent\textit{Editorial scope.} Counted unit: the Lemma whose source label is \texttt{RT} in the article (source0.tex lines 438--440, source label \texttt{RT}), delivered together with the article's own complete original proof of it (source0.tex lines 441--483, one \texttt{proof} environment of 43 lines). No mathematics is written, repaired or re-derived: statement and proof are the article's own text line for line. Required context retained uncounted: hL (lines 364--366). Every definition, display and label of those lines is retained unchanged. Omitted from this package and not redistributed: 1--363, 367--437, 484--1070. Nothing inside a retained excerpt is omitted or altered: every retained body is delivered exactly as the article's lines are written, so no omission receipt of any kind accompanies this package. Citations: the retained excerpts carry no citation command at all, so no bibliography entry of the article is transcribed and no reference is redistributed. Reference adaptation: none was needed and none was made; every reference inside the delivered unit resolves inside it, so no unresolved reference or citation is produced. Printed slips kept literal: no printed word, symbol, hypothesis or number of the counted statement or of its proof was altered. Layout and class: the article's own class is \texttt{article} with packages including geometry, diagbox, mathtools, bbm, latexsym, epsfig, amsmath, amsthm, amssymb, enumerate, nicefrac, xcolor, graphicx, color; the wrapper is the lane's pinned \texttt{amsart} editorial wrapper at 10pt on A4, which restates the theorem-like environments the retained text needs and adds the article's own macro definitions that the retained text needs (\texttt{\textbackslash Pr}, \texttt{\textbackslash bfrac}, \texttt{\textbackslash brac}, \texttt{\textbackslash d}, \texttt{\textbackslash m}, \texttt{\textbackslash t}, \texttt{\textbackslash wh}), with extra wrapper packages (bm, bbm, dsfont, xspace, cancel, mathtools). The delivered numbering is the unit's own. Assets: no retained span contains a figure environment, a graphics-inclusion command, a PDF-page inclusion, a table or a TikZ picture; no image file is copied into this package, the package declares no asset dependency and no image file is redistributed with it. Licence: version arXiv:2506.03375v1 of this article is distributed under the Creative Commons Attribution 4.0 International licence, as recorded in the arXiv licence metadata for this exact version and shown on the abstract page https://arxiv.org/abs/2506.03375v1. A full rights notice is delivered at licenses/arxiv-2506.03375-CC-BY-4.0.txt. Characters: every retained excerpt is pure ASCII, so the delivered engine, XeLaTeX with its default Latin Modern text font, reports no missing character.
\begin{equation}\label{hL}
h= \frac{d}{2 \log d}, \qquad L=  \frac {100d}{\log d}.
\end{equation}

\begin{lemma}\label{RT}
Let $p \ge p_c$, and let $X_t$ be a random walk on $Q_{n,p}$. Then w.h.p. for all $v \in V$ and all $T =O(\log^k n)$, $k $ constant, $R_v(T)=1+O(1/\log d)$.
\end{lemma}

\begin{proof}
As in \eqref{hL} of P3, fix the values of $h,L$ to $h= \frac{d}{2 \log d}$ and $ L=  \frac {100d}{\log d}$.
Let $t_0=d/\log^2 d$.
The proof is in three steps, from $t \le  t_0$, from $t_0<t \le h$, and from $h \le t \le T=\log^k n$.

We first consider the case where, with the possible exception of $v$ itself, no vertex within distance $h$ of $v$ has degree at most $L$.
\begin{align*}
\sum_{\t=1}^{t_0} \Pr(X_\t=v) \le& \sum_{\t=1}^{t_0} \frac 1L \Pr(X_{\t-1}\in N(v))
\le \frac{t_0}{L} = O \bfrac 1{\log d}.
\end{align*}
Let $v_i$ be a vertex at distance $1 \le i \le h$ from $v$. The probability the distance to $v$ decreases to $i-1$ at the
next step is at most $i/L$, and the probability it increases to $i+1$ is at least $(L-i)/L$. Thus the drift
away from $v$ per step is at least
\[
\mu \ge \frac{L-h}{L}- \frac hL = \frac{L-2h}{L} = \frac{99}{100}.
\]
In $t \le h$ steps the expected displacement of the walk from $v$ is at least $t\mu$. Let $dist(X_t,v)$ be the
actual displacement. For $\d>0$ constant
\[
\Pr(dist(X_t,v)\le (1-\d)t \mu) \le e^{-\Omega(\d^2 t\m)}.
\]
Thus
\[
\sum_{\t=t_0}^h \Pr(X_\t=v) \le he^{-\Omega(\d^2 t_0\m)} =o\bfrac{1}{\log d}.
\]
It follows  that, w.h.p., in $h$ steps the walk is at least distance $H=\mu h(1-2\d)$ from $v$,
where $H \ge 2 d/(5 \log d)$, say. Let
\[
\wh q= \frac{L-h}{L}= \frac{199}{200}, \qquad \wh p= \frac hL= \frac 1{200}.
\]
Consider a biassed random walk  with transition probabilities $\wh p$ of one step left and $\wh q$ of one step right,
setting out from $H-1$ on the integer line $\{0,1,...,H\}$. The probability the walk reaches the origin $v$ before returning to $H$ is
\[
\frac{ \bfrac{\wh q}{\wh p}-1}{ \bfrac{\wh q}{\wh p}^H-1}=O(\mu^{-h})=O\brac{ \bfrac{199}{200}^{2d/5\log d}}.
\]
Thus, with $T=\log^k n$, for any constant $k$,
\[
\sum_{\t=h}^{T} \Pr(X_\t=v) \le O\brac{ \bfrac{199}{200}^{2d/5\log d}}=o\bfrac{1}{\log d}.
\]

Next consider the case where  vertex $w$ is one of the at most 2 vertices of degree at most $L$ is within distance $h$ of $v$.
If $w \in N(v)$ this can increase the expected returns by $O(1/L)$. Suppose $w$ is a distance $i\ge 2$ from $v$. In the worst case assume the walk always returns to level $i-1$ (a wasted move). Deleting all edges between $w$ and its neighbours leaves all vertices within distance $h$ of $v$ with degree at least $L-1$. This has a negligible effect on the analysis given above.
\end{proof}

\end{document}
